% =====================================================
% 01_introduction.tex
% =====================================================
\todo{full draft — 1 page}

A trained generative adversarial network~\cite{Karras2019StyleGAN,
Karras2020StyleGAN2} defines a smooth map $G : \W \to \Img$ from a
Euclidean latent space to an ambient image space. Recent
interpretability work has produced a rich vocabulary for asking
questions about this map: which spatial regions does an attribute
boundary affect~\cite{Bau2019GANDissection, Shen2020InterFaceGAN}?
What unsupervised directions exist~\cite{Harkonen2020GANSpace,
Shen2021SeFa}? How disentangled are the resulting
edits~\cite{Wu2021StyleSpace, Eastwood2020DisentanglementMetrics}?

We argue that these questions admit a unified treatment through
\emph{differential geometry}. The generator's image set $\Mfd = G(\W)$
is a finite-dimensional immersed submanifold of $\R^{3HW}$, and the
above questions correspond, respectively, to (i)~the pushforward of
tangent vectors through $G$, (ii)~the spectrum of $dG$ restricted to
the manifold tangent space, (iii)~the structure of the
\emph{second}-order pushforward, and (iv)~the commutator (Lie bracket)
of attribute-induced vector fields.

This paper makes three contributions.

\textbf{1. A geometric framework.} We formalise six structural claims
about trained GAN manifolds (\S\ref{sec:theory}, claims C1--C6) that
together describe the tangent and curvature structure of $\Mfd$ and how
it interacts with attribute editing.

\textbf{2. An efficient computational realisation.} For our regime —
single-direction input, many-pixel output — \emph{forward-mode}
automatic differentiation (\jvp) is the right algorithmic direction.
A single composed \jvp computes first- and second-order per-pixel
sensitivities in $\sim 2.2 \times$ the cost of a vanilla generator
forward pass. We further exploit this for unsupervised attribute
discovery via tangent-space clustering.

\textbf{3. An applied payoff: compositional editing prediction.} We
show that the \emph{mixed Hessian} $d^2 G(b_a, b_b)$, also estimable
via composed \jvp, predicts when two attribute boundaries combine
non-linearly when applied simultaneously. The same machinery that
visualises individual attribute saliency tells us which edit
combinations are safe and which will interfere.

We validate every claim on three pre-trained generative domains
(LSUN bedroom, FFHQ faces, LSUN church) and compare against five
baselines spanning the segmentation-based dissection
\cite{Bau2019GANDissection}, unsupervised
factorisation~\cite{Harkonen2020GANSpace, Shen2021SeFa}, per-channel
style analysis~\cite{Wu2021StyleSpace}, and gradient attribution
\cite{Selvaraju2017GradCAM} families. A user study on local-edit
precision confirms that saliency-mask-guided editing is preferred over
global editing by human raters.

The framework is classifier-free: every signal comes from the
generator's own derivative structure, not from a separately trained
attribute classifier. This removes a major source of bias and
re-training overhead that has constrained prior GAN attribution
methods.

% ---- structure-of-paper paragraph ----
The remainder of the paper proceeds as follows.
\S\ref{sec:related} positions the work against
GAN-interpretation, classical saliency, Riemannian deep-learning
geometry, and forward-mode autodiff literatures.
\S\ref{sec:theory} introduces the geometric framework and the six
claims.
\S\ref{sec:method} describes the \jvp-based estimators.
\S\ref{sec:experiments} validates each claim across three domains
with quantitative metrics and baselines.
\S\ref{sec:application} shows the compositional-editing-prediction
result and a saliency-guided local editing user study.
\S\ref{sec:discussion} discusses limitations, particularly the
restriction to constant vector fields (linear boundaries) and
implications for non-linear edit methods like
EditGAN~\cite{Ling2021EditGAN}.
